\documentclass[11pt,letterpaper]{article}
\usepackage[margin=1in]{geometry}
\usepackage{amsmath,amssymb,amsthm}
\usepackage{enumitem}
\usepackage{booktabs}
\usepackage{xcolor}
\usepackage{tcolorbox}
\usepackage{microtype}
\tcbuselibrary{breakable,skins}

\newtcolorbox{objbox}[1]{breakable,enhanced,
  colback=red!3,colframe=red!60!black,
  fonttitle=\bfseries,title=Referee Objection #1,
  boxrule=0.8pt,left=5pt,right=5pt,top=4pt,bottom=4pt}

\newtcolorbox{respbox}{breakable,enhanced,
  colback=green!3,colframe=green!55!black,
  fonttitle=\bfseries,title=Response,
  boxrule=0.8pt,left=5pt,right=5pt,top=4pt,bottom=4pt}

\newcommand{\kt}{\tilde{\kappa}}
\newcommand{\kB}{k_{\mathrm{B}}}
\newcommand{\Scert}{S_{\mathrm{cert}}}
\newcommand{\Leff}{\Lambda_{\mathrm{eff}}}
\newcommand{\EP}{E_P}

\title{\textbf{Referee-Proofing Appendix}\\[4pt]
\large Objections and Responses Specific to the Operational Definition\\
and Cooper-Pair Constraint\\[3pt]
\normalsize Supplement to: \textit{Entanglement-Entropy Sourcing of Gravity} (v4)}

\author{Kevin Monette}
\date{March 2026}

\begin{document}
\maketitle

This appendix simulates a skeptical but intelligent referee focusing specifically on the
two most vulnerable aspects of the v4 framework: the operational definition of $\Scert$
and the Cooper-pair constraint.  For each objection, we state the manuscript modification
or computation needed to answer it fully.

% ─────────────────────────────────────────────────────────────────────────────
\section*{Objection 1}

\begin{tcolorbox}[breakable,enhanced,colback=red!3,colframe=red!60!black,
  fonttitle=\bfseries,title={Referee Objection 1: The definition is circular},
  boxrule=0.8pt]
``The paper defines $\Scert$ as entanglement generated by the protocol $P$, certified by
a witness $\mathcal{W}$, and varied between branches.  But this is circular: you are
defining your gravitational source as `whatever differs between your two branches.'
Nothing prevents the experimenter from choosing two branches that differ in energy,
temperature, electromagnetic fields, or any other quantity --- and calling that difference
$\Scert$.  How does the framework distinguish genuine entanglement-sourced gravity from
any other differential effect?''
\end{tcolorbox}

\begin{tcolorbox}[breakable,enhanced,colback=green!3,colframe=green!55!black,
  fonttitle=\bfseries,title={Response},boxrule=0.8pt]
The referee raises a legitimate concern that must be addressed explicitly in the manuscript.
The answer has two parts.

\textbf{Part 1: The definition is not circular because of condition (iii) (witness certification).}
$\Scert$ is not ``whatever differs between branches.'' It is specifically the entanglement
entropy of the half-system bipartition of the device state, certified by an entanglement
witness $\mathcal{W}$.  An entanglement witness is an observable $W$ satisfying
$\mathrm{Tr}(W\rho) \geq 0$ for all separable states and $\mathrm{Tr}(W\sigma) < 0$ for some
entangled state $\sigma$.  The condition $\mathrm{Tr}(W\rho^A) < 0$ certifies that the
state \emph{in branch A} is genuinely entangled (not classically correlated), independent
of branch B.  The differential $\Delta\Scert = \Scert^A - \Scert^B$ then measures the
difference in certified quantum entanglement, not any arbitrary property.

\textbf{Part 2: Matched energy and classical correlators are the primary systematics,
not part of the definition.}
The experimental protocol requires that branches A and B have identical single-particle
density matrices, identical two-body correlators, and identical energy expectation values.
Any residual differential that is not attributable to quantum entanglement (e.g., from
differential thermal dissipation, different mean photon number, or classical magnetic
correlations) is a systematic to be controlled, not a contribution to $\Scert$.

\textbf{Manuscript modification needed:} Add one paragraph to Section~2.2 making the
witness condition explicit and distinguishing $\Scert$ from generic branch differences.
A sentence such as: ``The witness requirement ensures that $\Scert > 0$ certifies genuine
quantum entanglement; any branch difference attributable to classical correlations,
differential energy, or systematic offsets is excluded by design and must be bounded
by the experimental systematic error budget.''
\end{tcolorbox}

% ─────────────────────────────────────────────────────────────────────────────
\section*{Objection 2}

\begin{tcolorbox}[breakable,enhanced,colback=red!3,colframe=red!60!black,
  fonttitle=\bfseries,title={Referee Objection 2: The Cancellation Proposition assumes what it needs to prove},
  boxrule=0.8pt]
``Proposition~6.1 assumes that the initial total state is a product state
$\rho_T^{(0)} = \rho_D^{(0)} \otimes \rho_{B_g}^{(0)}$.  But this is precisely what
is in question: in a superconducting device, the qubit degrees of freedom are physically
embedded in a superconducting substrate.  There is unavoidable coupling between the qubit
modes and the Cooper-pair condensate.  The assumption of a product initial state is not
justified; it is assumed for convenience and then the conclusion is that background
correlations cancel.  This is circular.''
\end{tcolorbox}

\begin{tcolorbox}[breakable,enhanced,colback=green!3,colframe=green!55!black,
  fonttitle=\bfseries,title={Response},boxrule=0.8pt]
This objection is technically correct and must be addressed quantitatively, not
just conceptually.

\textbf{The product-state assumption is an approximation, not an axiom, and the
approximation error is quantifiable.}  For superconducting transmon qubits, the
qubit-substrate coupling is dominated by phononic backaction with strength
$J_{\mathrm{ph}} \lesssim 1$\,kHz, compared to qubit transition frequencies
$\omega/2\pi \sim 5$\,GHz.  The ratio $J_{\mathrm{ph}}/\omega \sim 2\times10^{-7}$
quantifies how far the true initial state deviates from a product state.  The device-substrate
entanglement in the initial state is of order $J_{\mathrm{ph}}/\omega$ per qubit,
which for 50 qubits gives a background $\Scert \lesssim 50 \times 10^{-7} \approx 10^{-5}$\,bits.

For the Planck-scale benchmark ($E_* = E_P$, $\Delta\Scert^{\mathrm{device}} = 25$\,bits):
\[
\frac{\Delta\Scert^{\mathrm{background}}}{\Delta\Scert^{\mathrm{device}}}
\lesssim \frac{10^{-5}}{25} = 4\times10^{-7}.
\]
This ratio quantifies the error in the product-state approximation.  It is negligible
for the stated experimental purpose.

\textbf{For non-superconducting platforms} (neutral atoms, trapped ions), the
device-background coupling is even weaker, and the approximation is even better justified.

\textbf{Manuscript modification needed:} Add a quantitative estimate of the
product-state approximation error in the discussion following Proposition~6.1.
State explicitly that it is an approximation with quantifiable error
$\delta\Scert/\Scert^{\mathrm{device}} \sim J_{\mathrm{ph}}/\omega \sim 10^{-7}$,
not an assumption taken without justification.
\end{tcolorbox}

% ─────────────────────────────────────────────────────────────────────────────
\section*{Objection 3}

\begin{tcolorbox}[breakable,enhanced,colback=red!3,colframe=red!60!black,
  fonttitle=\bfseries,title={Referee Objection 3: OP-new is not a conjecture, it is an assumption},
  boxrule=0.8pt]
``The paper states Conjecture~6.2 (thermal state conjecture) as an open problem.  But
the \emph{entire experimental program} depends on this conjecture being true.  If thermal
equilibrium states have nonzero $\Scert$ under realistic operations, then every piece of
matter in the universe sources gravity entanglement-dependently, and the framework is
already ruled out by the most basic gravitational measurements.  Calling this a `conjecture'
understates how fundamental it is.  The paper should either prove it or acknowledge that
the framework is contingent on an unproven assumption that could falsify it entirely.''
\end{tcolorbox}

\begin{tcolorbox}[breakable,enhanced,colback=green!3,colframe=green!55!black,
  fonttitle=\bfseries,title={Response},boxrule=0.8pt]
The referee is correct that OP-new is more fundamental than the paper currently
acknowledges.  This requires a stronger statement in the manuscript.

\textbf{The correct framing is:}
\begin{enumerate}[noitemsep]
\item The Cancellation Proposition (Prop.~6.1) handles the \emph{differential} observable.
  It is a theorem, not a conjecture, and it does not depend on OP-new.
\item OP-new (Conjecture~6.2) concerns the \emph{absolute} gravitational source from
  equilibrium matter, which is a separate question.  For the \emph{experiment}, only
  the differential matters --- the Cancellation Proposition is sufficient.
\item For the \emph{theory's consistency with existing gravity measurements}, OP-new
  is required.  The framework would be inconsistent with existing data if thermal
  $\Scert$ is large under realistic operations.
\end{enumerate}

\textbf{The correct manuscript language} should say:
\begin{quote}
``For the experimental program proposed here, only the differential $\Delta\Scert$ matters;
Proposition~6.1 establishes that background contributions cancel exactly regardless of
whether Conjecture~6.2 holds.  However, consistency of the framework with existing
gravitational measurements (which show no anomalous contribution from room-temperature
matter) requires that either (a)~Conjecture~6.2 holds (thermal $\Scert = 0$) or
(b)~$E_* \ll E_P$ (making all signals undetectable).  The Planck-scale experimental
program is viable only if both (a) holds and the Cancellation Proposition applies.''
\end{quote}

\textbf{Partial proof of Conjecture~6.2:} For a toy model of $N$ non-interacting qubits
at temperature $T$ in a product state $\rho_\beta = \bigotimes_k \rho_k^\beta$, the
bipartite entanglement across any cut is exactly zero (product states are separable).
This establishes Conjecture~6.2 for the free-qubit model.  For interacting systems
(BCS), the SSR-accessible entanglement is bounded by the Wiseman-Vaccaro
formula~\cite{Wiseman2003} and can be estimated numerically for specific Hamiltonians.

\textbf{Manuscript modification needed:} Expand the OP-new discussion to (a)~clearly
separate the differential question (handled by Prop.~6.1) from the absolute question
(OP-new), (b)~state explicitly that the experimental program requires only Prop.~6.1,
and (c)~provide the free-qubit partial proof as supporting evidence.
\end{tcolorbox}

% ─────────────────────────────────────────────────────────────────────────────
\section*{Objection 4}

\begin{tcolorbox}[breakable,enhanced,colback=red!3,colframe=red!60!black,
  fonttitle=\bfseries,title={Referee Objection 4: $E_*$ is unfalsifiable},
  boxrule=0.8pt]
``The paper introduces $E_*$ as a free parameter and then observes that only $\Leff =
\kt E_* \alpha_{\mathrm{screen}}$ is constrained.  But this means a null experimental
result can always be ``explained'' by choosing a smaller $E_*$.  The theory can never be
falsified because for any experimental sensitivity, one can postulate $E_* \ll E_P$ and
the signal disappears.  This is not a falsifiable framework; it is a framework with
a permanent escape hatch.''
\end{tcolorbox}

\begin{tcolorbox}[breakable,enhanced,colback=green!3,colframe=green!55!black,
  fonttitle=\bfseries,title={Response},boxrule=0.8pt]
The referee is raising the most important strategic objection.  The response requires
distinguishing between falsifying \emph{the specific Planck-scale hypothesis} and
falsifying ``the framework'' as a whole.

\textbf{The Planck-scale hypothesis is falsifiable.}
The statement ``if $E_* = E_P$, $|\kt| = 1/4$, and $\alpha_{\mathrm{screen}} = 1$, then
$\Delta a = 3.6\times10^{-9}$\,m/s$^2$ for 100 bits at 100\,$\mu$m'' is a specific,
falsifiable numerical prediction.  A null result at this sensitivity with verified
entanglement would falsify this specific combination.

\textbf{The general framework is not falsifiable in isolation.}  This is true.  $E_*$
is a free parameter, and a null result always permits $E_* \to 0$.  The paper should
acknowledge this explicitly rather than claiming the framework as a whole is falsifiable.

\textbf{The correct falsification structure:}
\begin{center}\small
\begin{tabular}{@{}p{5.5cm}p{5cm}@{}}
\toprule
What is falsifiable & What is not \\
\midrule
$E_* = E_P$, $|\kt| = 1/4$, $\alpha_{\mathrm{screen}} = 1$ (a specific point in parameter space) &
The general framework for any $(E_*, \kt, \alpha_{\mathrm{screen}})$ \\
\addlinespace
$|\Leff| > 10^{-10}$\,J/bit (from a specific force-sensing experiment) &
Whether $E_*$ exists at all \\
\addlinespace
$\alpha_{\mathrm{screen}}$ for a specific 50-qubit circuit (from OP3 numerics) &
The coupling $\kt$ in the absence of a UV theory \\
\bottomrule
\end{tabular}
\end{center}

\textbf{Manuscript modification needed:} Add a clear paragraph in the Conclusions
stating: ``The framework as a whole, with $E_*$ as a free parameter, is not independently
falsifiable. What is falsifiable is the specific Planck-scale hypothesis $E_* \approx E_P$
combined with the holographic regulator estimate $\kt = -1/4$.  A null result at the
specified sensitivity would rule out this specific parameter combination, not the
existence of entanglement-gravity coupling in general.''
\end{tcolorbox}

% ─────────────────────────────────────────────────────────────────────────────
\section*{Objection 5}

\begin{tcolorbox}[breakable,enhanced,colback=red!3,colframe=red!60!black,
  fonttitle=\bfseries,title={Referee Objection 5: The superselection rule argument does not apply to neutral atoms},
  boxrule=0.8pt]
``The paper uses superselection-rule-based arguments (Wiseman-Vaccaro) to argue that
Cooper-pair entanglement is not operationally accessible.  But the proposed flagship
experiment uses neutral atoms in optical tweezers, not superconductors.  Neutral atoms
have no charge superselection rule; their hyperfine degrees of freedom are fully accessible.
How does the framework exclude background thermal entanglement in atomic gases, where
the accessible entanglement argument does not apply?''
\end{tcolorbox}

\begin{tcolorbox}[breakable,enhanced,colback=green!3,colframe=green!55!black,
  fonttitle=\bfseries,title={Response},boxrule=0.8pt]
The referee correctly identifies a gap in the argument.  For neutral atoms, the SSR
argument does not apply directly, and the Cooper-pair exclusion reasoning must be replaced
by a different one.

\textbf{The correct exclusion mechanism for neutral atoms is the Cancellation Proposition,
not the SSR argument.}  The Cancellation Proposition (Prop.~6.1) does not rely on
superselection rules.  It relies only on:
\begin{enumerate}[noitemsep]
\item The thermal background of the gas being identical in branches A and B.
\item The entangling protocol acting only on the designated qubit register (a $6\times6$
  array of Rb atoms in tweezers), not on the thermal gas surrounding the array.
\end{enumerate}
For a neutral-atom tweezer array, the thermal photon bath in the vacuum chamber and the
residual background gas are identical in both branches (both branches use the same trap,
same vacuum pressure, same temperature).  The Cancellation Proposition applies directly
without requiring a superselection rule.

\textbf{For the atomic gas itself:}  If the experiment uses a magneto-optical trap or
a BEC surrounding the tweezer array, the gas atoms are in a thermal state $\rho_\beta^{\mathrm{gas}}$
that is identical in branches A and B.  By Prop.~6.1, their contribution to
$\Delta\Scert$ is exactly zero.

\textbf{The SSR argument} (Wiseman-Vaccaro) provides additional support specifically for
superconducting platforms by showing that even the accessible entanglement of the Cooper-pair
background is small.  It is a \emph{bonus} for superconducting architectures, not a necessary
component of the general argument.

\textbf{Manuscript modification needed:} Clarify in the cancellation theorem discussion
that the SSR argument is a separate, platform-specific bonus for superconductors.  The
primary exclusion mechanism for all platforms is the branch-independence of background
correlations, established by Prop.~6.1.
\end{tcolorbox}

% ─────────────────────────────────────────────────────────────────────────────
\section*{Summary: Priority Manuscript Modifications}

\begin{center}\small
\begin{tabular}{@{}p{0.5cm}p{6cm}p{2.5cm}p{3cm}@{}}
\toprule
\# & Modification & Section & From objection \\
\midrule
1 & Add explicit witness-condition paragraph distinguishing $\Scert$ from generic branch differences & 2.2 & Obj.~1 \\
2 & Add quantitative product-state approximation error ($J_{\mathrm{ph}}/\omega \sim 10^{-7}$) after Prop.~6.1 & 6.2 & Obj.~2 \\
3 & Separate differential (Prop.~6.1, no conjecture needed) from absolute (OP-new required) questions & 6.3 & Obj.~3 \\
4 & Add paragraph stating Planck-scale hypothesis is falsifiable; general $E_*$ framework is not & Conclusions & Obj.~4 \\
5 & Clarify SSR argument is a superconductor bonus; Prop.~6.1 handles all platforms & 6.2 & Obj.~5 \\
\bottomrule
\end{tabular}
\end{center}

Modifications 1, 3, 4, and 5 are writing tasks (estimated 2 hours total).
Modification 2 requires one numerical estimate from the literature (IBM calibration data),
then one sentence.  All five modifications can be incorporated into a v4$\to$v4.1
revision without structural changes to the paper.

\begin{thebibliography}{9}
\bibitem{Wiseman2003}
H.~M.~Wiseman and J.~A.~Vaccaro,
\textit{Phys.\ Rev.\ Lett.}\ \textbf{91}, 097903 (2003).

\bibitem{Verstraete2003}
F.~Verstraete and J.~I.~Cirac,
\textit{Phys.\ Rev.\ Lett.}\ \textbf{91}, 010404 (2003).
\end{thebibliography}

\end{document}
